A grain only follows the membrane as long as the membrane does not accelerate downwards faster than $g$. At the top of its swing, the membrane accelerates downwards with $\hat a = A\,\omega^2$, while a grain on it falls only with $g$: once $\hat a > g$, the membrane drops away below it. So the grains lift off when
\begin{align}
 A\,\omega^2 &= g
\end{align}

\begin{abcliste}
\abc
\begin{align}
 \omega &= \omF = \sqrt{\frac{\ncg}{\A}} = \om
\end{align}
and with $\omega = 2\pi f$
\begin{align}
 f &= \fF \\
   &= \frac{\om}{2\pi} \\
   &= \f \approx \result{\fP{0}{2}}
\end{align}

\abc
\begin{align}
 A_2 &= \AbF \\
   &= \frac{\ncg}{\left(2\pi\times\fb\right)^2} \\
   &= \Ab \approx \result{\AbP{-3}{2}}
\end{align}
At higher frequencies a tiny amplitude is enough: that is why salt dances on loudspeakers already for deep bass notes.
\end{abcliste}
